% \newglossaryentry%
%     {entry-label}
%     {% key-value-list
%         % ::: REQUIRED ARGUMENTS :::
%         name = {item-name},%
%         description = {item-description},%
%
%         % ::: OPTIONAL ARGUMENTS :::
%         % Label (NOT name) of parent entry.
%         % Useful for hierarchical glossaries.
%         % Parent entry must be defined *before* this entry.
%         parent = {parent-entry-label},
%
%         % Plural form of description. 
%         % If omitted, `description` will be used.
%         descriptionplural = {item-description-plural-form},
%
%         % How entry will appear in text on subsequent use (2nd use and onward). 
%         % If omitted, `name` will be used.
%         text = {text-appearance},
%
%         % How entry will appear in text on first use.
%         % If omitted, `text` will be used.
%         first = {first-text-appearance},
%
%         % Plural form of entry in text
%         % If omitted, `text` will be used with `\glspluralsuffix` appended.
%         plural = {plural-text-appearance},
%
%         % Plural form of first entry appearance. 
%         % If omitted, value from `plural` is used.
%         firstplural = {first-plural-text-appearance},
%
%         % Symbol associated with with entry.
%         % Not all styles will display symbol!
%         symbol = {item-symbol},
%
%         % Same as symbol, but for plural.
%         % If omitted, value is set to `symbol` key.
%         symbolplural = {item-symbol-plural},
%
%         % Text to be used by sort comparator when ordering glossary.
%         % If omitted, `name` field will be used, unless package option 
%         % `sort` is set to 'def' or 'use'. See `glossaries-user.pdf:130 (141)`
%         sort = {sort-text},
%
%         % Label of glossary, in which this entry belongs.
%         % If omitted, `\glsdefaulttype` is used.
%         type = {glossary-label}
%
%         % For further key options, see `glossaries-user.pdf` on pages:
%         %   - p. 133 (144)
%         %   - p. 368 (379)
%         %   - p. 376 (387)
%     }

% \newglossaryentry{relu}%
% {%
%     name = {ReLU},
%     description = {Rectified Linear Unit activation function}
% }

\newacronym{wysiwyg}{WYSIWYG}{What You See Is What You Get}
\newacronym{ctan}{CTAN}{Comprehensive \TeX\ Archive Network}
\newacronym{doi}{DOI}{Digital Object Identifier}
\newacronym{mwe}{MWE}{minimal working example}


% For long descriptions you may use:
% \longnewglossaryentry%
%     {entry-label}
%     {% key-value-list
%         name = {item-name},  %   REQUIRED!!
%
%         % For possible keys, see above
%     }
%     {description} %             REQUIRED!!

% For acronyms you may use:
% \newacronym%
%     [key=value] % Optional key-value-list. Appended to \newglossaryentry kv-list
%     {acronym-label}
%     {acronym-abbreviation}
%     {acronym-long-form}



